\section{Height-one characters and first-level elliptic torsion}\label{sec:first-level-refinements}
This section proves the all-prime height-one assertion, exact Hesse defects over both a specified field and a variable-ramification family, and the residue-field obstruction over a fixed cyclotomic base.

\subsection{The finite scalar in height one}
The proof identifies the finite descent scalar directly from the first torsion layer.

\begin{proof}[Proof of Proposition~\ref{prop:height-one-index}]
A height-one group in the allowed filtration class has a single nonzero graded piece, so itself acquires good reduction. Its dimension $d$ is zero or one. In dimension zero the good model is \'etale; in dimension one its Cartier dual is \'etale. Over the inertia of a good-reduction field its Tate character is therefore respectively $1$ or $\chi_p$. On the original inertia group it has the form
\[
                         \rho=\chi_p^d\epsilon,
\]
where $\epsilon$ is a finite $\Z_p^\times$-valued character. Indeed $\rho/\chi_p^d$ is trivial on the finite-index inertia subgroup of a good field. This same finite scalar is the special descent character. For the \'etale model it is the scalar on the constant Tate module. For the multiplicative model the assertion follows by Cartier duality; the extra generic cyclotomic factor contributes no additional special descent scalar.

Suppose first that $p$ is odd. Lemma~\ref{lem:principal-congruence-torsion} shows that $1+p\Z_p$ is torsion-free, so reduction embeds every finite subgroup of $\Z_p^\times$ into $\F_p^\times$. Since $\mu_p\subset K$, one has $\chi_p\equiv1\pmod p$. Thus the actual representation $M[p]$ determines $\epsilon\bmod p$, and its unique Teichm\"uller lift recovers every value of $\epsilon$. Recovering the dimension is unnecessary for the special-character conclusion.

For $p=2$, the finite subgroup of $\Z_2^\times$ is $\{1,-1\}$. To see this, reduction modulo four has image of order two, while $1+4\Z_2$ is torsion-free: a nonzero leading term $2^a b$ with $a\ge2$ gains exactly one valuation under squaring and retains its valuation under odd powers. Every odd-order element of the finite inertia image therefore has special scalar one. The requested prime-to-two character is constant one and is determined by $M[2]$. The index is nonnegative by definition, so its least value is zero in every prime.
\end{proof}

\subsection{An explicit Hesse specialization}
The following calculation proves both full torsion for the particular equation and the exact ramification needed for its good reduction.

\begin{proof}[Proof of Theorem~\ref{thm:explicit-hesse-twelve}]
The polynomial $X^2+1$ is separable and irreducible modulo three, so $\Q_3(i)/\Q_3$ is unramified quadratic. Over that field $\pi=\zeta-1$ satisfies the Eisenstein polynomial $X^2+3X+3$. Thus $e_K=2$, the residue field is $\F_9$, and $v_K(\pi)=1$. In particular
\[
 \pi^2=-3(\pi+1),\qquad \pi^3=6\pi+9,\qquad \pi^6=-27.
\]
First allow $t$ to be a variable and put
\[
 A=-27t(t^3+8),\qquad B=54(t^6-20t^3-8).
\]
Direct multiplication gives
\begin{equation}\label{eq:hesse-invariants}
 \Delta=2^{12}3^9(t^3-1)^3,\qquad c_4=2^4 3^4t(t^3+8).
\end{equation}
For the first identity, put $U=t^3$ and use
\[
 (U^2-20U-8)^2-U(U+8)^3=-64(U-1)^3.
\]
The equation is therefore smooth whenever $t^3\ne1$.

Let $\tau=\zeta-\zeta^2$, so $\tau^2=-3$. Its eight nonzero $3$-torsion points are
\begin{align*}
 x_*&=-9t^2,& y_*&=\pm12\tau(t^3-1),\\
 x_j&=3t^2+12(\zeta^jt+\zeta^{-j}),&
 y_j&=\pm36(\zeta^jt^2+\zeta^{-j}t+1),\qquad j=0,1,2.
\end{align*}
Here is a direct verification. The three-division polynomial is
\[
 \psi_3(x)=3x^4+6Ax^2+12Bx-A^2.
\]
It follows from the tangent slope $(3x^2+A)/(2y)$ and the equality $x(2P)=x(P)$; for an affine point with $y\ne0$, that equality gives $2P=-P$. Polynomial division yields
\[
 \frac{\psi_3(x)}{3(x+9t^2)}
 =x^3-9t^2x^2+27t(t^3-16)x-27(t^3+8)^2.
\]
After $x=3t^2+12w$, the cubic on the right becomes
\[
 1728\bigl(w^3-3tw-(t^3+1)\bigr),
\]
whose roots are $w=\zeta^jt+\zeta^{-j}$. At $x_*$ the cubic equation has right side $-432(t^3-1)^2=y_*^2$. At $x_0=3(t+2)^2$ it has right side $1296(t^2+t+1)^2$. For general $j$, replace $t$ by $t'=\zeta^{-j}t$ and multiply $x_0(t')$ by $\zeta^{-j}$; the equation is carried back to the original one and gives the displayed $y_j$.

The four $x$-coordinates are distinct: an equality $x_j=x_k$ for $j\ne k$ forces $t=\zeta^{-j-k}$, while $x_j=x_*$ forces $t^2+\zeta^jt+\zeta^{-j}=0$. Either condition implies $t^3=1$. A displayed $y$ can vanish only under the same condition. Thus the eight points are distinct nonidentity points killed by three, and exhaust the nonzero geometric $3$-torsion in characteristic zero. This proves full rational torsion for the actual curve over any field containing $\mu_3$ with $t^3\ne1$.

Now specialize to $t=1+i\pi$. Exact reduction gives
\begin{align*}
 t^3&=10-9i+(9-3i)\pi,\\
 t^3-1&=9(1-i)+3(3-i)\pi,\\
 t^3+8&=9(2-i)+3(3-i)\pi.
\end{align*}
In each of the last two lines the summands have valuations four and three. Hence $v_K(t^3-1)=v_K(t^3+8)=3$, and \eqref{eq:hesse-invariants} gives $v_K(\Delta)=27$ and $v_K(c_4)=11$.
The change $x=\pi^4X$, $y=\pi^6Y$ gives the exact integral model
\begin{equation}\label{eq:hesse-minimal-model}
 Y^2=X^3+aX+b,\qquad
 a=-3-12i-8i\pi,\quad b=-30+12i-20\pi.
\end{equation}
For example,
\[
 t^6-20t^3-8=-27(15-6i+10\pi)=\pi^6(15-6i+10\pi),
\]
which verifies the constant coefficient after scaling; the linear coefficient follows from $-27=\pi^6$. Both $a$ and $b$ have valuation one. The new discriminant has valuation three, so the model is minimal: an integral model with smaller nonnegative discriminant valuation would differ by a positive multiple of twelve. Its special fibre is the cusp $Y^2=X^3$, hence reduction is additive. Its $c_4$ has valuation three and $j$ has valuation six.

Let $r$ be a root of $X^3+aX+b$. This polynomial is Eisenstein, so $K(r)/K$ is totally ramified of degree three and $v_K(r)=1/3$. The element $c=a+3r^2$ has valuation one, since $v_K(3r^2)=8/3$. Choose $u$ with $u^4=c$. In $K(r)$ the normalized valuation of $c$ is three, so adjoining this fourth root has exact degree and ramification index four. Thus $L=K(r,u)$ has degree and ramification index twelve over $K$, with unchanged residue field.
The substitution $X=r+u^2Z$, $Y=u^3W$ gives
\[
                     W^2=Z^3+(3r/u^2)Z^2+Z.
\]
The middle coefficient has valuation $2+1/3-1/2=11/6>0$. Its special fibre $W^2=Z^3+Z$ has discriminant $-64\ne0$ and three-division polynomial $-1$ in characteristic three. It is smooth and supersingular, proving actual good reduction over $L$.

For the matching lower bound, let $F/K$ give good reduction. The discriminant transformation forces $12\mid3e(F/K)$, hence $4\mid e(F/K)$. Since two is invertible, the $2$-torsion of a good elliptic scheme is finite \'etale, and all its points become rational over an unramified extension of $F$. The point $(r,0)$ in \eqref{eq:hesse-minimal-model} has $v_K(r)=1/3$. Unramified extensions preserve the value group, so $3\mid e(F/K)$. Every good-reduction index is divisible by twelve, and $L$ attains it. These arguments remain valid after unramified base change. Finally $v(j)>0$ excludes multiplicative reduction after extension, since a nodal minimal fibre has $c_4$ a unit and $\Delta$ of positive valuation. Thus the exact semistability defect is twelve.
\end{proof}

\subsection{A variable-ramification Hesse family}
The same Hesse equation realizes exact wild defects over an infinite range of complete bases.

\begin{proof}[Proof of Theorem~\ref{thm:hesse-variable-defect}]
Put $s=\varpi^r$, $t=1+s$, $S=e+r$ and $\delta=t^3-1$. Since $2r<e$,
\[
 v_K(\delta)=v_K(t^3+8)=3r.
\]
Indeed the terms in $\delta=3s+3s^2+s^3$ have values $e+r,e+2r,3r$, and the last is strictly smallest; adding nine does not change that value. With $A=-27t(t^3+8)$ and $B=54(t^6-20t^3-8)$, the identity
\[
 t^6-20t^3-8=\delta^2-18\delta-27
\]
gives $v_K(B)=3e+6r$. The other two terms have larger values because $2r<e$. Thus
\begin{equation}\label{eq:variable-hesse-values}
 \begin{aligned}
 v_K(A)&=3S,&v_K(B)&=3e+6r,\\
 v_K(\Delta)&=9S,&v_K(c_4)&=4e+3r,&v_K(j)&=3e>0.
 \end{aligned}
\end{equation}
Here the invariant identities are~\eqref{eq:hesse-invariants}. In particular $\delta\ne0$. The eight rational points and their universal division-polynomial verification in the proof of Theorem~\ref{thm:explicit-hesse-twelve} apply to every such $t$ over a field containing $\mu_3$. They therefore prove full rational $3$-torsion for this actual curve.

We find the cubic factor in a good-reduction field using the simpler polynomial
\[
                         q(z)=z^3-3tz+2.
\]
The fractional-linear expression
\begin{equation}\label{eq:hesse-two-torsion-map}
              \alpha(z)=\frac{3((t^3-4)z-6t^2)}{2+tz}
\end{equation}
satisfies the exact polynomial identity
\[
 \bigl(\alpha(z)^3+A\alpha(z)+B\bigr)(2+tz)^3
                       =-1728\delta^2q(z).
\]
Its determinant is $24\delta\ne0$, and $q(-2/t)=8(1-t^{-3})\ne0$. It therefore bijects the roots of $q$ with the roots of the two-division cubic and gives $K(\alpha)=K(z)$ for corresponding roots. The polynomial $q$ also results by setting $X=Y$ in the Hesse plane cubic $X^3+Y^3+Z^3-3tXYZ=0$ and putting $z=Z/X$. The displayed fractional-linear identity supplies the field comparison directly.

Set $z=1+w$. Then
\[
                     q(1+w)=w^3+3w^2-3sw-3s.
\]
The coefficient points $(1,S)$ and $(2,e)$ lie strictly above the line joining $(0,S)$ to $(3,0)$: $S>2S/3$ and $e>S/3$, the latter following from $r<2e$. Hence every root has value $v_K(w)=S/3$. This also follows directly because either a larger or a smaller value would give a unique least-valued term. Since $3\nmid S$, the denominator forces the cubic to be irreducible and totally ramified of degree three. Put $K_1=K(z)=K(\alpha)$.

The original cubic $X^3+AX+B$ has a single Newton segment: its coefficient at degree one lies strictly above the segment from degree zero to degree three, since
\[
                    3e+3r>\frac23(3e+6r)=2e+4r.
\]
Thus $v_K(\alpha)=e+2r$. For $c=A+3\alpha^2$ the summands have values $3e+3r$ and $3e+4r$, so $v_K(c)=3S$ and $v_{K_1}(c)=9S$ in the normalized integer valuation. Choose $u$ with $u^4=c$. Every unit of $K_1$ has a fourth root by Hensel's lemma, because its residue field is algebraically closed of characteristic three. Removing that unit reduces the extension to a fourth root of a uniformizer power. Its degree and ramification index are exactly
\[
                     h=\frac4{\gcd(4,9S)}=\frac4{\gcd(4,S)}.
\]
Consequently $L=K_1(u)$ has degree and ramification index $3h$ over $K$.

The substitution $x=\alpha+u^2X$, $y=u^3Y$ gives
\[
                      Y^2=X^3+(3\alpha/u^2)X^2+X.
\]
Its middle coefficient has value
\[
                    e+(e+2r)-\frac{3S}{2}=\frac S2>0.
\]
The reduction is the smooth cubic $Y^2=X^3+X$ in characteristic three. Its three-division polynomial is $-1$, so it has no nonidentity geometric three-torsion and is supersingular. Hence $L$ gives actual good reduction at ramification degree $12/\gcd(4,S)$.

For the matching lower bound let $F/K$ be any finite good-reduction extension and put $d=e(F/K)$. Its valuation ring is strictly henselian, so the finite \'etale group scheme $E[2]$ of the good model is constant. In particular it contains a two-torsion coordinate $\alpha$, and hence $z$ by~\eqref{eq:hesse-two-torsion-map}. The value $v_K(z-1)=S/3$ forces $3\mid d$. A good Weierstrass equation has unit discriminant, and changes of Weierstrass equation multiply discriminants by twelfth powers. Therefore
\[
                         12\mid9dS,
\]
so $4/\gcd(4,S)\mid d$. This divisor is prime to three; their product divides every possible $d$, and $L$ attains it. Finally $v_K(j)>0$ excludes multiplicative reduction after any extension. The curve is potentially good by the displayed model, so the good-reduction and semistability defects coincide. This proves the claimed exact defect.
\end{proof}

\subsection{The first-level defect table}
The complete existence criteria supply actual elliptic curves at every index in the table.

\begin{proof}[Proof of Corollary~\ref{cor:three-first-level-table}]
The Weil pairing forces $\mu_3\subset K$, hence $e$ is even. Theorem~\ref{thm:ordinary-unramified-p} excludes an ordinary good special fibre for an additive curve with full rational $3$-torsion. Over an algebraic closure, a supersingular special fibre has equation $y^2=x^3-x$. Its origin-preserving automorphisms are
\[
 (x,y)\longmapsto(a^2x+r,a^3y),\qquad a^4=1,\quad r^3-r=0.
\]
Comparison of the coefficients in the general Weierstrass coordinate change gives these conditions, so the group has order twelve. Inertia embeds in it by smooth proper base change and faithfulness on a prime-to-three Tate module. The only possible nontrivial orders are therefore $2,3,4,6,12$; trivial inertia would give good reduction.

At $n=1$, Theorem~\ref{thm:tame-full-level} for $m=2,4$ gives exactly $me\ge8$ with positive even $e$. For $m=3,6,12$, Theorem~\ref{thm:three-range} has $s\ge0$ and exponent $\max(0,-s)=0$, giving the same inequality. The five rows follow immediately. Both theorems include actual every-index realizations with algebraically closed residue field, so this argument supplies existence as well as necessity. Proposition~\ref{prop:explicit-cubic-defect-six} gives the asserted example over the particular field $\Q_3(\zeta_3)$.
\end{proof}

\subsection{The residue obstruction at the critical index}
At equality in the basic ramification bound, normalized torsion coordinates must fit into the residue field of the original base.

\begin{lemma}\label{lem:critical-torsion-residue}
Let $E/K$ have full rational $p$-torsion and potentially supersingular reduction. Suppose it has good reduction over a finite Galois extension $L/K$ of ramification degree $m$. Then
\[
                            p^2-1\le me_K.
\]
If equality holds, there is an injective $\F_p$-linear map
\[
                            E[p](K)\longrightarrow k_K.
\]
In particular equality is impossible when $k_K=\F_p$.
\end{lemma}
The inequality is the first-level form of the full-division-field bound; in the number-field setting it is covered by \citet[Theorem~1.1, p.~2]{Smith2023}. The proof below includes the complete-field setting and the equality descent to $k_K$.
\begin{proof}
Choose an integral formal parameter $T$ at the origin of the good model over $\Ocal_L$, with $w(L^\times)=\Z$. Every $p$-torsion point reduces to the origin, because the good special fibre is supersingular. Its $p^2-1$ nonzero coordinates lie in the maximal ideal of $\Ocal_L$. The series $[p](T)/T$ has constant term $p$ and distinguished degree $p^2-1$. Weierstrass preparation gives a unit times a monic polynomial whose roots are exactly these coordinates. Its constant term gives
\[
                    \sum_{P\ne0}w(T(P))=w(p)=me_K.
\]
Each summand is a positive integer, proving the inequality.

At equality every summand equals one. Fix $P_0\ne0$, put $z=T(P_0)$ and define
\[
 \lambda(P)=\overline{T(P)/z},\qquad\lambda(0)=0.
\]
The integral formal law is $S+T$ plus terms of degree at least two, so $\lambda$ is additive. It is injective because every nonzero coordinate has valuation one. Hence it is an injective $\F_p$-linear map into $k_L$.

We must descend the image to $k_K$. Generic descent isomorphisms extend uniquely to the good models by the N\'eron mapping property \citep[Theorem~17.1 and Remark~17.2, p.~70]{Milne2008AV}. For $\sigma\in\Gal(L/K)$ their formal expressions are
\[
 F_\sigma(Z)=a_\sigma Z+\text{terms of degree at least two},
 \qquad a_\sigma\in\Ocal_L^\times.
\]
A rational point satisfies $T(P)=F_\sigma(\sigma T(P))$. Modulo valuation two, numerator and denominator of $T(P)/z$ transform by the same factor $a_\sigma$. Their residue ratio is fixed by every $\sigma$. Since $L/K$ is Galois and the residue field is perfect, the induced residue group is $\Gal(k_L/k_K)$, so $\lambda(P)\in k_K$. This proves the stronger equality statement for the actual finite torsion group.
\end{proof}

\begin{proof}[Proof of Corollary~\ref{cor:cyclotomic-full-p}]
Let $p\ge5$, $K=\Q_p(\zeta_p)$ and suppose such an elliptic curve exists. The shifted cyclotomic polynomial gives $e_K=p-1$ and $k_K=\F_p$. Theorem~\ref{thm:ordinary-unramified-p} excludes ordinary potential good reduction, so the good special fibre is supersingular.

For a prime $\ell\ne p$, potential good reduction makes the inertia image on $T_\ell E$ finite. Smooth proper base change identifies it with a faithful subgroup of the good special fibre's automorphism group. In characteristic at least five, a short equation $y^2=x^3+Ax+B$ shows that origin-preserving automorphisms are $(x,y)\mapsto(u^2x,u^3y)$, with $u^4A=A$ and $u^6B=B$. The translation and extra $y$-terms vanish because three and two are invertible. Thus this automorphism group is cyclic of order dividing four or six, and the finite inertia image has order $m\le6$.

Choose a sufficiently large $r$ that this finite inertia image remains faithful on $E[\ell^r]$. Then $L=K(E[\ell^r])$ is finite Galois with ramification degree exactly $m$. Its inertia acts trivially on the whole $T_\ell E$, so NOS gives good reduction over $L$ \citep[Theorem~1, p.~493]{SerreTate1968}. Lemma~\ref{lem:critical-torsion-residue} yields
\[
                         p^2-1\le m(p-1)\le6(p-1).
\]
This implies $p+1\le6$, excluding $p\ge7$. For $p=5$ it forces $m=6$ and equality. The lemma would then inject the $25$-element group $E[5](K)$ into $\F_5$, a contradiction. Finally Proposition~\ref{prop:explicit-cubic-defect-six} gives the example for $p=3$.
\end{proof}
